\documentclass[11pt]{article}
\usepackage[T1]{fontenc}
\usepackage{lmodern,amsmath,amssymb,microtype}
\usepackage[margin=1.02in]{geometry}
\usepackage[colorlinks=true,urlcolor=blue,citecolor=blue]{hyperref}
\setlength{\emergencystretch}{2em}
\hypersetup{pdfauthor={Wentao Li},pdftitle={Minimal domination in hereditary classes of incomparability graphs}}
\title{Minimal domination in hereditary classes of incomparability graphs}
\author{Wentao Li\\[2pt]\small Prepared with assistance from OpenAI Codex}
\date{November 2026}
\begin{document}
\maketitle
\begin{abstract}
For every fixed co-bipartite graph $H$, the inclusion-minimal dominating sets of an $H$-free incomparability graph can be listed in incremental polynomial time. This establishes Conjecture~7.3 of Bonamy, Defrain, Micek and Nourine. We give an explicit bound on the VC dimension of closed neighbourhoods and describe the resulting classification of hereditary incomparability classes.
\end{abstract}

\paragraph{A shattering lemma.} All graphs and orders in this note are finite. In an incomparability graph two different vertices are adjacent exactly when incomparable in an underlying partial order. A set $S$ is shattered by closed neighbourhoods if every $T\subseteq S$ is $N[v]\cap S$ for some vertex $v$. Suppose $H$ has a partition into two cliques of sizes $a,b\ge1$. Set $m=\max(1,\lceil\log_2(b^2)\rceil)$ and $k=2a+2m$. If $G$ is an incomparability graph with a shattered set of size at least $2k$, it contains an induced $H$.

Indeed, any three elements $l<m'<r$ in $S$ would rule out the trace containing $l,r$ and omitting $m'$. A witness outside this chain comparable to $m'$ is comparable to at least one of $l,r$ by transitivity, and a chain vertex cannot realize the trace either. Thus $S$ has height at most two. Its minimal elements and its nonminimal elements are two antichains, so it contains an antichain $A$ of size $k$. Subsets of shattered sets are shattered.

Divide $A$ into two disjoint blocks of $a$ vertices, each numbered by the first clique of $H$, and $m$ disjoint marker pairs. Let $P_i$ be the cross-neighbours of vertex $i$ of the second clique among the vertices of the first. For each $(i,j)\in[b]^2$ prescribe a trace equal to $P_i$ on the first block and $P_j$ on the second. Give these $b^2$ indices distinct $m$-bit identifiers. In each marker pair put precisely the vertex indicated by that bit into the trace. Choose witnesses $w_{ij}$. Their marker traces are distinct and have equal size, hence no two traces are nested. Each witness lies outside $A$, since a vertex of $A$ sees all of $A$ in its closed neighbourhood.

Each witness is comparable to some element of $A$ since it misses a marker. All its comparisons with $A$ have the same direction, because $A$ is an antichain. Colour the witness positive or negative according to this direction. Two distinct witnesses of the same colour cannot be comparable: for comparable witnesses of that direction their comparable-element sets in $A$ would be nested by transitivity, whereas their traces, and hence their complements, are incomparable on the markers. If every row $i$ of the witness array has a positive entry, pick one in every row. These give a clique realizing $H$ with the first block. Otherwise some row consists entirely of negative entries. Its entries give a clique realizing $H$ with the second block. This proves the lemma. Single-vertex or empty obstructions are immediate.

\paragraph{Sharp order of magnitude.} Take $1\le a\le b$. In the worst case over the forbidden graph with these side sizes the bound is $\Theta(a+\log(b+1))$ as $a+b$ grows. For the lower bound already take $H=K_a\sqcup K_b$. Use a co-bipartite host with sides $D$ of size $d$ and $2^D$, making cross-edges record membership. It is an incomparability graph: orient the edges of the bipartite complement uniformly from one side to the other to obtain a height-two order. The set $D$ is shattered. Taking $d=a-1$ excludes $H$, since neither clique of the host can intersect both anticomplete pieces of an induced $H$, and each piece has at least $a$ vertices. Separately take the largest $d$ with $d+2^d<a+b$. The host then has too few vertices for $H$, and $d\ge \log_2(b+1)-O(1)$. Taking the larger example proves the stated lower order.

\paragraph{Enumeration theorem.} The hypergraph $\mathcal N$ with vertex set $V(G)$ and edges $N[v]$ has as its inclusion-minimal transversals precisely the minimal dominating sets of $G$. Constructing it takes polynomial time. Deleting duplicate edges and edges properly containing other edges preserves its minimal transversals and cannot increase VC dimension. The lemma bounds the VC dimension by $2k-1$ on the entire promised input class. Mary's bounded-VC-dimension transversal enumeration theorem now gives an incremental polynomial algorithm: the running time for the next output or certification of completion is polynomial in the input and output obtained so far, for fixed $H$. In particular this implies the output-polynomial bound asked for in Conjecture~7.3. A transitive orientation of the complement need not be supplied to the enumeration algorithm.\hfill$\square$

\paragraph{Hereditary classification.} If a hereditary subclass of incomparability graphs omits some co-bipartite graph, apply the preceding theorem. If it contains every co-bipartite graph, general minimal-transversal enumeration reduces to domination on this class, by the co-bipartite reduction of Kant\'e, Limouzy, Mary and Nourine. The opposite reduction uses $\mathcal N$. Thus the second case is output-polynomial equivalent to unrestricted hypergraph transversal enumeration, whose general complexity remains open. Kobayashi, Kurita and Wasa recently proved that general transversal enumeration admits no output-polynomial algorithm under the Exponential Time Hypothesis (ETH), \url{https://arxiv.org/abs/2610.11473}. Consequently, under ETH the hereditary subclasses admitting output-polynomial domination enumeration are exactly those omitting some co-bipartite graph. The positive algorithm is unconditional.

\paragraph{The nonquantitative shortcut and priority.} Lozin's Theorem~12 already characterizes the minimal hereditary classes of unbounded neighbourhood VC dimension as bipartite, co-bipartite and split. Incomparability graphs omit the bipartite six-cycle and the split three-sun. Indeed an independent triple is a chain $l<m<r$, and there is no path from $l$ to $r$ avoiding $N[m]$: vertices strictly below and strictly above $m$ cannot be adjacent. The alternating vertices of the six-cycle and the three independent tips of the three-sun are asteroidal triples, in which each pair has such a path avoiding the third closed neighbourhood. Thus Lozin and Mary together also imply the qualitative result. The shattering lemma above gives the sharp worst-case order in terms of both clique sizes. We claim neither of the general tools as new.

\paragraph{Sources.} M. Bonamy, O. Defrain, P. Micek and L. Nourine, §7, \url{https://arxiv.org/abs/2004.07214} (2020 preprint, \href{https://doi.org/10.1016/j.disc.2025.114904}{2026 journal publication}). A. Mary, Theorem~1 and Corollary~12, \url{https://arxiv.org/abs/2407.00694}. V. Lozin, \emph{Hereditary classes of graphs: a parametric approach}, Discrete Applied Mathematics 325 (2023), Theorem~12. M. M. Kant\'e et al., \emph{On the enumeration of minimal dominating sets and related notions}, SIAM J. Discrete Math. 28 (2014), 1916--1929. Independent proof and priority review are pending.
\end{document}
